%%==================================================================
\section{What the method does not give}\label{sec:scope}
%%==================================================================

The summation routine needs three things of a factor system: scopes of size
$\lambda^{O(1)}$, occurrence numbers of size $\lambda^{O(1)}$ after the
branching of Lemma~\ref{lem:frequent}, and total degree $O(n)$. The following
problems fail one of them, and this paper proves nothing about them.

\begin{enumerate}[label=(\alph*),leftmargin=2.2em]
\item \emph{Dense objectives.} Max-Cut on a graph with $\Theta(n^2)$ edges has
total degree $\Theta(n^2)$ and $\Theta(n)$ occurrences per variable; branching
on frequent variables fixes $\Theta(n)$ of them. Theorem~\ref{thm:density}
reaches average degree $o(\sqrt{\log n/\log\log n})$, and by
Remark~\ref{rem:densityend} no design for the routine reaches total degree
$\frac14(n+1)\log n$. Corollary~\ref{cor:complement} covers the dense graphs whose
complement is sparse, and no other dense graphs. The
chromatic and Tutte polynomials as polynomials in $q$ need connectivity, which
is not a bounded sum of local factors; Theorem~\ref{thm:tutte} gives their
values at each fixed positive integer $q$ on bounded-degree graphs, and nothing
off the hyperbolas $H_q$. Section~\ref{ssec:indep} shows that the independence
polynomial escapes the density limit because a vertex of large degree forces
many others; no such forcing is available for a cut.
\item \emph{Problems in $\mathrm{P}$ with quadratic conjectured complexity.}
The lower bounds for Orthogonal Vectors, $3$SUM, All-Pairs Shortest Paths and
their relatives that were derived from the strong exponential time hypothesis
lose their foundation, as \oai{} note; but the reductions that gave them run
from CNF-SAT on $n$ variables to instances of size $2^{n/2}$, and an algorithm
with exponent $o(n)$ for the former gives nothing below $N^{2-o(1)}$ for the
latter.
\item \emph{$\mathrm{P}$ versus $\mathrm{NP}$.} A polynomial-time reduction
from an arbitrary problem in $\mathrm{NP}$ to $3$-SAT produces formulas with
$n^{C}$ variables for some $C>1$, and $2^{o(n^{C})}$ is not polynomial.
Nothing in this paper bears on the question.
\end{enumerate}

Each of the results of Sections \ref{sec:applications}, \ref{sec:more} and \ref{sec:snp} has
been stated with the exact hypothesis under which it is proved. The constants
in the exponents are explicit through Lemma~\ref{lem:params},
Theorem~\ref{thm:bounded-occ} and Theorem~\ref{thm:explicit}, and the dependence on the weight bound $\Omega$
is polynomial in every statement.
